% Aula de catalogação: registros bibliográficos em MARC 21
% Modelo do ÉguaTex (https://eguatex.github.io), em domínio público pela licença CC0 1.0:
% use, modifique e distribua sem pedir permissão e sem citar a origem.
%
% Compile com pdfLaTeX (o ÉguaTex faz isso sozinho).
%
% Slides com código ou registros MARC precisam da opção [fragile] no frame,
% porque o ambiente verbatim lê o texto exatamente como está escrito.

\documentclass[11pt]{beamer}

% ---------- idioma e fontes ----------
\usepackage[T1]{fontenc}
\usepackage[utf8]{inputenc}
\usepackage[brazilian]{babel}
\usepackage{lmodern}
\usepackage{tgschola}
\usepackage[varqu]{zi4}  % Inconsolata: fonte monoespaçada legível para os registros
\usefonttheme{serif}

% ---------- tema ----------
\usetheme{Boadilla}
\usecolortheme{beaver}
\setbeamertemplate{navigation symbols}{}
\setbeamertemplate{blocks}[rounded][shadow=false]

\usepackage{booktabs}
\usepackage{fancyvrb}

% Caixa para registros MARC: fundo claro e fonte monoespaçada pequena.
% Use \begin{registro} ... \end{registro} dentro de um frame [fragile].
\definecolor{fundomarc}{RGB}{246,241,232}
\newenvironment{registro}[1]
  {\setbeamercolor{block body}{bg=fundomarc, fg=black}%
   \begin{block}{#1}\footnotesize}
  {\end{block}}

% ---------- instruções de preenchimento ----------
% As notas em cinza explicam o que pôr em cada slide.
% Troque \instrucoestrue por \instrucoesfalse para esconder todas de uma vez.
\newif\ifinstrucoes
\instrucoestrue
\newcommand{\instrucao}[1]{\ifinstrucoes\par\vfill{\color{gray}\scriptsize\itshape #1\par}\fi}

% ---------- dados da capa: preencha aqui ----------
\title[Catalogação em MARC 21]{Catalogação descritiva em MARC 21}
\subtitle{Aula 5: campos, indicadores e subcampos}
\author[Nome do professor]{Nome do professor ou da professora}
\institute[UF Exemplo]{Representação Descritiva II\\Curso de Biblioteconomia, Universidade Federal Exemplo}
\date[Ano]{Cidade, Ano}

\begin{document}

\begin{frame}
  \titlepage
\end{frame}

% ======================================================================
\begin{frame}{Objetivos da aula}
  Ao final da aula, você deve ser capaz de:
  \begin{enumerate}
    \item reconhecer as três partes de um registro MARC;
    \item ler as etiquetas, os indicadores e os códigos de subcampo;
    \item codificar a descrição de um livro nos campos 100, 245, 264 e 300.
  \end{enumerate}
  \instrucao{Comece com até 4 objetivos escritos com verbos observáveis. Volte a eles no último slide.}
\end{frame}

% ======================================================================
\section{Estrutura}
% ======================================================================
\begin{frame}{As partes de um registro}
  \begin{table}
    \centering\small
    \begin{tabular}{@{}llp{5.2cm}@{}}
      \toprule
      Parte & Tamanho & Função \\
      \midrule
      Líder & 24 posições & dados para o sistema processar o registro \\
      Diretório & variável & mapa com etiqueta, tamanho e posição de cada campo \\
      Campos & variável & controle (00X) e dados (010 a 9XX) \\
      \bottomrule
    \end{tabular}
  \end{table}
  \begin{block}{Para lembrar}
    O catalogador preenche os campos. O líder e o diretório são gerados pelo sistema.
  \end{block}
  \instrucao{Use tabelas para comparar partes ou campos. Uma tabela com mais de 5 linhas fica ilegível no projetor.}
\end{frame}

\begin{frame}[fragile]{Anatomia de um campo}
  \begin{registro}{Campo 245: título e indicação de responsabilidade}
\begin{verbatim}
245 10 $a Título do livro de exemplo :
       $b subtítulo da obra /
       $c Nome Sobrenome.
\end{verbatim}
  \end{registro}
  \begin{itemize}
    \item \texttt{245}: etiqueta de três dígitos.
    \item \texttt{1} e \texttt{0}: indicadores (há entrada secundária de título; nenhum caractere a ignorar).
    \item \texttt{\$a}, \texttt{\$b}, \texttt{\$c}: códigos de subcampo.
  \end{itemize}
  \instrucao{Destaque um campo por slide. Explique cada parte logo abaixo do registro, na ordem em que aparece.}
\end{frame}

% ======================================================================
\section{Registro completo}
% ======================================================================
\begin{frame}[fragile]{Um livro descrito em MARC 21}
  \begin{registro}{Registro bibliográfico (exemplo)}
\begin{verbatim}
020    $a 9780000000000
040    $a XXXX $b por $c XXXX
082 04 $a 025.32 $2 23
100 1  $a Sobrenome, Nome, $e autor.
245 10 $a Título do livro de exemplo :
       $b subtítulo da obra / $c Nome Sobrenome.
250    $a 2. ed.
264  1 $a Cidade : $b Editora, $c Ano.
300    $a 000 p. : $b il. ; $c 23 cm.
650  4 $a Catalogação descritiva.
\end{verbatim}
  \end{registro}
  \instrucao{Monte o registro de um livro do acervo da disciplina. Troque os dados de exemplo e mantenha o alinhamento das colunas com espaços.}
\end{frame}

\begin{frame}{Campos mais usados em livros}
  \begin{columns}[T]
    \begin{column}{0.48\textwidth}
      \begin{block}{Descrição}
        \begin{description}[\texttt{264}]
          \item[\texttt{245}] título
          \item[\texttt{250}] edição
          \item[\texttt{264}] publicação
          \item[\texttt{300}] descrição física
        \end{description}
      \end{block}
    \end{column}
    \begin{column}{0.48\textwidth}
      \begin{block}{Pontos de acesso}
        \begin{description}[\texttt{264}]
          \item[\texttt{100}] autor pessoal
          \item[\texttt{110}] entidade
          \item[\texttt{650}] assunto
          \item[\texttt{700}] coautor
        \end{description}
      \end{block}
    \end{column}
  \end{columns}
  \instrucao{Agrupe os campos por função. Os alunos guardam melhor quatro campos por grupo do que uma lista corrida.}
\end{frame}

% ======================================================================
\section{Prática}
% ======================================================================
\begin{frame}[fragile]{Exercício}
  Codifique os dados abaixo nos campos 100, 245 e 264.
  \begin{registro}{Dados da folha de rosto (exemplo)}
\begin{verbatim}
Autora: Nome Sobrenome
Título: Manual de catalogação: teoria e prática
Local e editora: Cidade, Editora
Ano: Ano
\end{verbatim}
  \end{registro}
  \pause
  \begin{alertblock}{Dica}
    O subtítulo vai no \texttt{\$b}, depois de espaço, dois-pontos e espaço no fim do \texttt{\$a}.
  \end{alertblock}
  \instrucao{Dê um exercício curto para resolver em sala. A dica aparece no clique seguinte, com o comando \textbackslash pause.}
\end{frame}

\begin{frame}{Para estudar}
  \begin{itemize}
    \item Documentação oficial do formato bibliográfico MARC 21, mantida pela Library of Congress:\\
      \url{https://www.loc.gov/marc/bibliographic/}
    \item Manual de catalogação adotado pela disciplina (indique autor, título e edição).
  \end{itemize}
  \bigskip
  \begin{block}{Próxima aula}
    Campos de assunto (6XX) e vocabulários controlados.
  \end{block}
  \instrucao{Indique onde o aluno encontra a norma completa e anuncie o tema da próxima aula.}
\end{frame}

\end{document}
